\begin{abstract}
We disprove Hadwiger's conjecture by constructing arbitrarily large
graphs whose chromatic number exceeds their Hadwiger number. The examples
have independence number at most two, and even their ordinary fractional
chromatic number exceeds their Hadwiger number. Thus they also disprove
the fractional-coloring weakening discussed by Reed and Seymour.
\section{Introduction}\label{sec:introduction}

For a finite graph \(G\), write \(\alpha(G)\), \(\chi(G)\), and \(h(G)\)
for its independence number, chromatic number, and Hadwiger number.
Thus \(h(G)\) is the largest \(t\) for which \(G\) contains a \(K_t\)
minor.

Hadwiger's conjecture, formulated in 1943~\citep{hadwiger1943},
asserts that \(h(G)\ge\chi(G)\) for every finite nonempty simple
graph \(G\). We disprove it using graphs of independence number
at most two.

Two disjoint edges are \emph{touching} if an edge joins their
endpoint sets. A \emph{connected matching} is a matching whose edges
are pairwise touching; write \(\cm(G)\) for its maximum size. The word
connected in this definition requires adjacency of every pair of
matching edges.

Our main result is the following.
\begin{theorem}\label{thm:main}
There are arbitrarily large integers \(m\) for which an \(m\)-vertex
finite simple graph \(G\) satisfies
\[
             \alpha(G)\le2
             \qquad\text{and}\qquad
             \cm(G)<\frac{m}{100}.
\]
\end{theorem}

A complete minor with many singleton or two-vertex branch sets gives
a large connected matching. The precise counting inequality, proved in
\Cref{prop:matching-minor}, is
\[
                    h(G)\le\frac{|V(G)|+4\cm(G)+2}{3}.
\]
Every color class has size at most \(\alpha(G)\), so
\(\chi(G)\ge |V(G)|/\alpha(G)\). The same lower bound holds
for fractional coloring, giving a second consequence of the construction.
For a finite nonempty simple graph \(G\), let \(\mathcal I(G)\) be
the family of its independent vertex sets. A fractional coloring
assigns a nonnegative real weight \(w_I\) to every \(I\in\mathcal I(G)\)
such that \(\sum_{I\ni v}w_I\ge1\) for every vertex \(v\).
The ordinary fractional chromatic number \(\chi_f(G)\) is the minimum
of \(\sum_Iw_I\) over these assignments.
An ordinary coloring gives a fractional coloring with unit weights
on its color classes, so \(\chi_f(G)\le\chi(G)\).

\begin{corollary}\label{cor:hadwiger}\label{cor:fractional-hadwiger}
There are finite nonempty simple graphs \(G\) of arbitrarily large
order \(m\) for which
\[
 h(G)<\frac{26m}{75}+\frac23
 <\frac m2\le\chi_f(G)\le\chi(G).
\]
Thus Hadwiger's conjecture and its fractional-coloring weakening
\(\chi_f(G)\le h(G)\) are false.
\end{corollary}
\begin{proof}
Summing the vertex constraints of any fractional coloring gives
\[
